\chapter{Future Work}
\label{cap:futuro}

The following lines derive from the limitations in chapter~\ref{cap:discusion}.

\begin{enumerate}
  \item \textbf{Fee levels, latency and external volatility.} The simulator should incorporate the swap with fee of equation~\eqref{eq:comision}, a delay in the execution of arbitrage and a reference with volatility, in order to turn into results the factors that today are design candidates (table~\ref{tab:factores}). With a fee, the no-arbitrage band ceases to be a free parameter and is determined by $\gamma=-\log(1-f)$.
  \item \textbf{CEX with depth and latency.} Replacing the exogenous reference with an order book that responds to the AMM's trades would make it possible to study bidirectional feedback.
  \item \textbf{Concentrated liquidity.} Extending the AMM to positions over price intervals, which change the relation between reserve and impact.
  \item \textbf{Alternative provider rules.} Comparing the one-shot withdrawal rule with an optimal-stopping one \cite{bergault2025exit} and with gradual withdrawals, and repeating H2 with and without the panic agent in a pre-registered manner.
  \item \textbf{Design with adequate precision.} Preparing the pilot with the same shock design as the confirmation, so that the sample-size guidance is representative.
  \item \textbf{Reduced map.} Testing maps with more variables or with nonlinear terms, and evaluating whether they improve on the scalar autoregression; exploring other criteria for classifying regimes.
  \item \textbf{External validation.} Contrasting the error contraction rate and the relation between depth and impact with data from a pair traded on a CEX and a DEX.
  \item \textbf{Pedagogical evaluation.} Measuring with students whether the interactive interface improves the understanding of the mechanisms.
\end{enumerate}
